\chapter{Zadatak 1 --- što je točno zadano}
\label{ch:z1-zadatak}

\section{Izvorni tekst, doslovno}
\label{sec:z1-izvorni}

Ovo je stranica koju je mentor potpisao. Ništa u sivom okviru nije dodano,
prepričano ni protumačeno.

\begin{zadatak}{Zadatak 1 (QP, konveksnost, KKT, dualnost)}
\skenzad[0.99]{z1-original}
\end{zadatak}

\subsection{Prijepis, znak po znak}

Zbog čitljivosti isti tekst slijedi i pretipkan. Formulacija je nepromijenjena.

\begin{zadatak}{Zadatak 1 --- prijepis}
\textbf{Zadatak 1. (QP, konveksnost, KKT, dualnost)} Zadan je optimizacijski problem
\begin{align}
&\min_{x} \quad x_1^2 + 2x_2^2 - 0.3x_1x_2 + 2x_1 - 3x_2 \tag{1a} \\
&\text{uz ograničenja} \quad a_{11}x_1 + a_{12}x_2 \leq -10, \tag{1b} \\
&\phantom{\text{uz ograničenja}} \quad a_{21}x_1 + a_{22}x_2 \leq 3, \tag{1c}
\end{align}
gdje su numeričke vrijednosti koeficijenata
$a_{11} = -1$, $a_{12} = -1$, $a_{21} = 1$, $a_{22} = -1$.

\begin{enumerate}[label=\alph*),leftmargin=*]
  \item Je li ovo konveksan optimizacijski problem? Obrazložiti odgovor.
  \item Postavite KKT uvjete optimalnosti i riješite ih (odredite vrijednosti
        projektnih varijabli i Lagrangeovih množitelja koji zadovoljavaju te uvjete).
  \item Nacrtajte dozvoljeni skup, nivo krivulje funkcije cilja te izračunajte i
        skicirati gradijente ograničenja i funkcije cilja u točki koja zadovoljava
        KKT uvjete.
  \item Riješite problem u MATLABu (quadprog i/ili yalmip) --- odredite primarne i
        dualne varijable u optimumu te ih usporedite s rezultatima (Lagrangeovim
        množiteljima) iz (b).
  \item Postavite i riješite Lagrangeov dualan problem (koristite MATLAB za
        numeričko rješavanje dualnog problema). Usporedite vrijednost izračunatih
        dualnih varijabli s onima dobivenim pod (d) i (b).
  \item Što možete reći o osjetljivosti optimalne vrijednosti ($p^\star =$ opt.
        vrijednost primarnog problema) na perturbacije u ograničenjima (na
        perturbacije koeficijenata $b_1$ i $b_2$)?
\end{enumerate}
\end{zadatak}

\section{Što u zadatku \emph{ne} piše}
\label{sec:z1-nije-receno}

Ovo je najvažnija stranica cijelog poglavlja. Sve što slijedi ispod nije dio
zadatka --- to su odluke i tumačenja koja mora donijeti onaj tko zadatak rješava.
Na obrani se očekuje da ih znaš navesti \emph{kao svoje odluke}, a ne da ih
predstaviš kao da su zadane.

\begin{longtable}{@{}p{0.36\linewidth} p{0.60\linewidth}@{}}
\toprule
\textbf{Što nije rečeno} & \textbf{Kako je protumačeno i zašto} \\
\midrule
\endfirsthead
\toprule
\textbf{Što nije rečeno} & \textbf{Kako je protumačeno i zašto} \\
\midrule
\endhead
Što su $b_1$ i $b_2$ iz podzadatka f). & U tekstu zadatka slovo $b$ se prije toga
uopće ne pojavljuje. Iz konteksta su to desne strane ograničenja (1b) i (1c),
dakle $b_1 = -10$ i $b_2 = 3$. Perturbacija ,,koeficijenata $b_1$ i $b_2$``
znači pomicanje tih desnih strana. \\[6pt]
U kojem obliku pisati ograničenja. & Zadatak ih daje kao $a\trans x \leq b$.
Notacija s predavanja i vježbi je $Gx - h \leq 0$, dakle sve na lijevu stranu i
nula desno. U ovom udžbeniku i u seminaru koristi se profesorova notacija, jer
KKT uvjeti u njoj imaju najjednostavniji oblik. \\[6pt]
Treba li koristiti \texttt{quadprog} ili \texttt{yalmip}. & Zadatak kaže
,,quadprog i/ili yalmip``, dakle bilo koje. U seminaru je odabran
\texttt{quadprog}, jer izravno vraća i Lagrangeove množitelje u polju
\texttt{lambda.ineqlin}, što je upravo ono što podzadatak d) traži usporediti. \\[6pt]
Koliko decimala navesti. & Nije zadano. Rješenje je egzaktno racionalno
($190/33$ i $140/33$), pa se u seminaru navodi razlomak, a decimale se pišu na
šest mjesta uz njega. \\[6pt]
Treba li dokazati da minimum postoji. & Ne traži se izrijekom, ali se
podrazumijeva pod ,,obrazložiti odgovor`` u a). Bez dokaza egzistencije tvrdnja
da je nađena točka \emph{globalni} minimum nije potpuna. \\[6pt]
Treba li provjeriti Slaterov uvjet. & Ne spominje se. Ipak je nužan: bez njega
KKT uvjeti su samo nužni, ne i dovoljni, pa se ne bi smjelo tvrditi da je
KKT točka optimum. U seminaru se provjerava točkom $x = (7,\,6)$. \\
\bottomrule
\end{longtable}

\begin{zamka}[Zamka: ,,koraci rješavanja`` koji nisu profesorovi]
U repozitoriju uz zadatak postoji i prijepis \texttt{TASK/zadatak\_omrius.md} u
kojem uz svaki podzadatak stoji numerirani popis \emph{Koraka rješavanja}
(,,zapisati Lagrangian``, ,,razmotriti sve kombinacije aktivnih ograničenja`` i
slično). Taj popis \textbf{nije} dio zadatka --- napisan je naknadno kao pomoć.
Ako se na obrani pozoveš na njega kao na tekst zadatka, izlažeš se pitanju koje
ne možeš dobro odgovoriti. Sve što je zadano stane u sivi okvir iz
\S\ref{sec:z1-izvorni}.
\end{zamka}

\section{Karta poglavlja}
\label{sec:z1-karta}

\begin{longtable}{@{}p{0.07\linewidth} p{0.42\linewidth} p{0.45\linewidth}@{}}
\toprule
& \textbf{Podzadatak} & \textbf{Gdje je obrađen} \\
\midrule
\endhead
a) & Konveksnost & Poglavlje~\ref{ch:z1-konveksnost} \\
b) & KKT uvjeti i njihovo rješavanje & Poglavlje~\ref{ch:z1-kkt} \\
c) & Grafički prikaz & Poglavlje~\ref{ch:z1-graficki} \\
d) & Numeričko rješenje (\texttt{quadprog}) & Poglavlje~\ref{ch:z1-matlab} \\
e) & Lagrangeov dualni problem & Poglavlje~\ref{ch:z1-dual} \\
f) & Osjetljivost optimalne vrijednosti & Poglavlje~\ref{ch:z1-osjetljivost} \\
--- & Pitanja, zamke i modelski odgovori & Poglavlje~\ref{ch:z1-obrana} \\
\bottomrule
\end{longtable}

Prije toga, u Poglavlju~\ref{ch:z1-formulacija}, zadatak se prevodi iz riječi u
matrični zapis --- korak koji zadatak ne traži izrijekom, a bez kojega se
nijedan od podzadataka ne može uredno riješiti.
